\chapter{Hereditary Discrepancy and Factorization Norms}
\label{ch:gamma2}

We have now seen several tools for proving upper bounds on the
discrepancy of set systems and matrices, some of which also give
efficient algorithms for computing a coloring achieving the
discrepancy bound. In this chapter we address the question of the
complexity of approximating discrepancy and hereditary discrepancy. In
addition to showing some hardness results, we will show that
hereditary discrepancy is approximated well by the
$\gamma_2$ matrix factorization norm. Beyond giving an approximation
algorithm for hereditary discrepancy, this result has a number of
applications to discrepancy theory. It gives a powerful tool for
giving short proofs of upper and lower bounds on discrepancy, and for
proving various stability properties of discrepancy. 

\section{Complexity of Computing Discrepancy}

Perhaps the most basic computational question about discrepancy is
whether $\disc(A)$ can be computed efficiently given a matrix $A$ as
input. An easier question is whether the promise that $\disc(A)$ is
small, \eg, at most a universal constant, or even $0$, is enough to
efficiently find a coloring whose discrepancy is smaller than
Spencer's worst-case bound $O(\sqrt{n\log(2m/n)})$. Unfortunately, the
answer to both these questions is negative, as long as $\mathsf{P}
\neq \mathsf{NP}$, as shown in the following theorem.

\begin{theorem}\label{thm:ch7-disc-hard}
  For a set system $(\SS, [n])$ with $|\SS| = O(n)$ sets, it is $\mathsf{NP}$-hard to distinguish between the cases: (i) $\disc(\SS) = 0$,
or (ii) $\disc(\SS) = \Omega(\sqrt{n})$.
% \begin{enumerate}
% \item $\disc(\SS) = 0$;
% \item $\disc(\SS) = \Omega(\sqrt{n})$.
% \end{enumerate}
\end{theorem}

We will establish Theorem~\ref{thm:ch7-disc-hard} by reduction from
the $2$-$2$ Set Splitting problem, whose definition and hardness are given below. %The result is due to Guruswami~\cite{venkat-SS}.

\begin{theorem}\label{thm:ch7-SS}(Guruswami~\cite{venkat-SS})
For a set system $(\TT, [n])$ with $m = O(n)$ sets, where $|T| = 4$
for all $T \in \TT$, and each $j \in [n]$ appears in at most $4$
sets, it is $\mathsf{NP}$-hard to distinguish between the cases:
\begin{enumerate}
\item $\disc{(\TT)} = 0$.
\item for all $\col \in \{\pm 1\}^n$, $|\{T \in \TT: \col(T) \not=
  0 \}| \geq \alpha m$, where $\alpha \approx 1/22$. %is a small enough constant. 
\end{enumerate}
\end{theorem}

Note that Theorem~\ref{thm:ch7-SS} already shows that it is
$\mathsf{NP}$-hard to decide if a set system has discrepancy $0$, or
discrepancy at least $2$. By composing a $2$-$2$ Set Splitting
instance, in an appropriate way, with a set system satisfying a robust
discrepancy lower bound of $\Omega(\sqrt{n})$, we can amplify this gap and prove Theorem~\ref{thm:ch7-disc-hard}. The incidence matrix
of the set system used in this composition is given by the following lemma. In it, we
use the notion of a Hadamard matrix, i.e., an $m \times m$ matrix with
entries in $\{-1,+1\}$ such that $H^T H = mI$. Such matrices are known
to exist, and can be efficiently constructed, for all $m$ that are
powers of $2$ (among other values of $m$). 

\begin{lemma}\label{lm:ch7-hadamard}
  Let $H$ be an $m\times m$ Hadamard matrix such that
  one of its rows is the all-ones vector. Then the matrix
  $A$ obtained by replacing each $-1$ entry of
  $H$  with $0$ satisfies
  \(
  \|A\col\|_\infty \ge \frac{1}{3}\|\col\|_2
  \)
  for  any $\col \in \R^m$.
\end{lemma}
\begin{proof}
  Let $\col \in \R^m$ be arbitrary. We first observe that, by
  comparing the $\ell_\infty$ and $\ell_2$ norms, 
  \begin{align}
  \label{eq:hadamard}
    \| H\col \|_{\infty}^2
    \geq \frac{1}{m} \|H\col \|_2^2
    = \frac1m \col^T  H^T H \col
    = \col^T \col =\|\col \|_2^2.
  \end{align}
  Notice that $A = \frac12 (H + J)$ where $J$ is the $m\times m$
  all-ones matrix. Suppose, without loss of generality, that the first
  row of $H$ is the all-ones vector, and, therefore, the first row of
  $A$ is also the all-ones vector. Then
  \[
    \|A\col\|_\infty = \frac12 \|H\col + J\col\|_\infty
    = \frac12 \max_{i = 1}^m |(H\col)_i + (A\col)_1|.
  \]
  By \eqref{eq:hadamard}, we know that $|(H\col)_{i^*}|
  \ge \|\col\|_2$ for some $i^*$. If $|(A\col)_1| \le \frac13\|\col\|_2$, then, by the
  triangle inequality,
  \[
    \|A\col\|_\infty %\ge %\frac12 |(H\col)_{i^*} + (A\col)_1|
    \ge \frac12 (|(H\col)_{i^*}| - |(A\col)_1|) \ge  \frac13 \|\col\|_2.
  \]
  Otherwise, if $|(A\col)_1| > \frac13\|\col\|_2$, then 
  \(
  \|A\col\|_\infty  \ge |(A\col)_1| > \frac13\|\col\|_2.
  \)
 % The lemma holds in both cases.
\end{proof}

Notice that the requirement that the Hadamard matrix $H$ has an all-ones
row is not restrictive. We can replace every entry $H(i,j)$ in $H$ by
$H(i,j)H(1,j)$, and still get a Hadamard matrix, but now all entries
in the first row are equal to $1$. 

Theorem~\ref{thm:ch7-SS} and Lemma~\ref{lm:ch7-hadamard} allow for a
very simple hardness reduction that shows a statement analogous to
Theorem~\ref{thm:ch7-disc-hard}, but for matrices with entries in
$\{0,4\}$, as opposed to set systems. The reduction simply takes the
$m\times n$ incidence matrix $B$ of the set system $\TT$ in
Theorem~\ref{thm:ch7-SS}, pads it with $0$'s if necessary to make sure
that $m$ is a power of $2$, and outputs a matrix $\mat \eqdef A B$ for
the $m\times m$ matrix $A$ defined in Lemma~\ref{lm:ch7-hadamard}.

\begin{exercise}
  Analyze the reduction described in the paragraph above.
\end{exercise}

In order to get a reduction that, instead, outputs a set
system, we need a simple decomposition lemma, stated next. 

\begin{lemma}\label{lm:ch7-partition}
  Suppose that $(\TT, [n])$ is a set system such that each set in $\TT$
  has at most $s$ elements, and each element $j \in [n]$
  appears in at most $t$ sets in $\TT$. Then $\TT$ can be partitioned
  in polynomial time
  into $k \le st+1$ disjoint set systems $\TT_1, \ldots, \TT_k$, where
  each $\TT_i$ consists of disjoint subsets of $[n]$.
\end{lemma}
\begin{proof}
  We partition $\TT$ greedily, similarly to the greedy coloring
  algorithm for bounded degree graphs. We list the sets in $\TT$ in some
  order as $T_1, \ldots, T_m$, and, for each $i$, starting
  from $i=1$ up to $i=m$, we add $T_i$ into $\TT_\ell$, for the smallest $\ell$
  for which $\TT_\ell$ does not contain any set intersecting
  $T_i$. Since each $T_i$ can intersect at most $st$ other sets in
  $\TT$, there is always some $\ell \le st+1$ so that $T_i$ can
  be added to $\TT_\ell$.
\end{proof}

\begin{proof}[Proof of Theorem~\ref{thm:ch7-disc-hard}]
  Let $\TT$ be the set system from Theorem~\ref{thm:ch7-SS}, and let
  us partition it into $k\le 17$ set systems $\TT_1, \ldots, \TT_k$,
  each consisting of disjoint sets, using
  Lemma~\ref{lm:ch7-partition}. Let $B_\ell$ be the incidence matrix
  of $\TT_\ell$. We can assume that each $B_\ell$ has dimensions $m$
  by $n$ for some $m$ which is a power of $2$, by padding with
  zeros. Let us take $H$ to be an $m\times m$ Hadamard matrix, and
  define $A$ as in Lemma~\ref{lm:ch7-hadamard}. We now define $\mat_\ell \eqdef A B_\ell$, for
  $\ell \in [k]$. Notice that, since $\TT_\ell$ consists of disjoint
  sets, $B_\ell$ has at most a single $1$ in each column, so
  $A B_\ell$ is an $m\times n$ matrix with entries in $\{0,1\}$. We
  can then interpret $\mat_\ell$ as the incidence matrix of a set
  system $\SS_\ell$ on $[n]$. We define the set system $\SS$ as the
  union $\SS_1\cup \ldots \cup \SS_k$. $\SS$ has at most $km \le 17m =
  O(n)$ sets and can clearly be constructed in polynomial time.

  To finish the proof, we claim that the following two implications
  hold for an absolute constant $c > 0$:
  \begin{align}
    \disc(\TT) = 0 &\implies \disc(\SS) = 0;\label{eq:ch7-compl-sets}\\
    \min_{\col \in \{-1,+1\}^n} |\{T \in \TT: \col(T) \not=  0 \}| \geq \alpha m
                   &\implies \disc(\SS) \ge c\sqrt{\alpha m}.\label{eq:ch7-sound-sets}
  \end{align}
  The theorem then follows from Theorem~\ref{thm:ch7-SS}.

  We first prove \eqref{eq:ch7-compl-sets}. Since $\disc(\TT)$ is $0$,
  there exists a coloring $\col \in \{-1,+1\}^n$ for which $\disc(\TT,
  \col) = 0$. Then $B_\ell \col = 0$ for all $\ell \le k$, since each
  row of each $B_\ell$ is the indicator vector of a set in $\TT$. Therefore,
  \[
    \disc(\SS,\col) = \max_{\ell = 1}^k \disc(\SS_\ell,\col)
    =
    \max_{\ell=1}^k \|\mat_\ell \col\|_\infty
    =
    \max_{\ell=1}^k \|AB_\ell \col\|_\infty = 0.
  \]

  Finally, we prove \eqref{eq:ch7-sound-sets}. Take any $\col \in
  \{-1,+1\}^n$, and suppose that $|\{T \in \TT: \col(T) \not=  0 \}|
  \geq \alpha m$. This means that for some $\ell \le k$,
  \[
    |\{T \in \TT_\ell: \col(T) \not=  0 \}| \geq \alpha m/k
    \ge \alpha m/17.
  \]
  Therefore, $B_\ell\col$ has at least $\frac{\alpha m}{17}$ non-zero
  entries. Since all entries in $B_\ell\col$ are integers, this means
  that $\|B_\ell\col\|_2 \ge \sqrt{\alpha
      m/17}$. So, by Lemma~\ref{lm:ch7-hadamard},
  \[
    \disc(\SS,\col) \ge \|\mat_\ell\col\|_\infty
    =
    \|AB_\ell \col\|_\infty
    \ge \frac13 \sqrt{\alpha m/17}. 
  \]
  Since $\col \in \{-1,+1\}^n$ was arbitrary, this completes the proof
  of \eqref{eq:ch7-sound-sets}, 
%  (with $c = \frac{1}{3\sqrt{17}}$), 
and,
  therefore, the theorem. 
\end{proof}

\section{Approximating Hereditary Discrepancy}
  
In contrast with the strong hardness result for discrepancy given by
Theorem~\ref{thm:ch7-disc-hard}, we will see in the rest of this
chapter that hereditary discrepancy can be approximated efficiently up
to multiplicative poly-logarithmic factors. This is one concrete
example of hereditary discrepancy being a ``better behaved''
discrepancy notion. We will see other examples later on, also in this
chapter.

\subsection{Upper Bounds}

Recall that the hereditary discrepancy $\herdisc(\mat)$ of a matrix
$\mat$ is the maximum discrepancy $\disc(\mat(*,J))$ for any subset
$J$ of the columns of $\mat$. Since this is a maximum over an exponential
number of quantities, each of them intractable to approximate, it is
far from obvious that we can find useful tractable upper bounds on
$\herdisc(\mat)$. In this subsection we show that this is, in fact,
possible. 

Our starting point is the observation that if $f$ is a function on
$m\times n$ matrices which is monotonically non-increasing under
taking submatrices, \ie, $f(\mat(*,J)) \le f(\mat)$ for any $J
\subseteq [n]$, and that also satisfies $\disc(\mat) \le f(\mat)$, then we
automatically get the stronger conclusion $\herdisc(\mat) \le
f(\mat)$. Indeed, for any $J \subseteq [n]$, we have
\[
  \disc(\mat(*,J)) \le f(\mat(*,J)) \le f(\mat).
\]
We call such a function $f$ a {\em hereditary upper bound} on discrepancy.

We use the notation $\|\mat\|_{1\to 2}$ for the maximum $\ell_2$ norm
of a column of the matrix $\mat$, and $\|\mat\|_{2\to\infty}$ for the
maximum $\ell_2$ norm of a row of the matrix $\mat$.\footnote{More
  generally, for an $m\times n$ matrix $\mat$, and
  $p,q \in [1,\infty]$,
  $\|\mat\|_{p\to q} \eqdef \sup_{x\neq 0} \frac{\|\mat
    x\|_q}{\|x\|_p}$. }

Notice first that both $c\sqrt{\log(2m)}\|\mat\|_{1\to 2}$ and
$c\sqrt{\log(2m)}\|\mat\|_{2\to\infty}$ are hereditary upper bounds on
discrepancy (for some constant $c>0$). The first follows from the Banaszczyk  upper bound for
the Koml\'os problem (Theorem~\ref{thm:komlos-bana}), and the second
from Hoeffding's inequality and a union bound, applied to a uniformly
random coloring.
Both these upper bounds can be, unfortunately, very loose, {\em e.g.},~consider
the all-ones matrix. 

However, we already saw in Lemma~\ref{lm:fact-ub}
that these bounds can be interpolated in a useful way using
Banaszczyk's theorem. Lemma~\ref{lm:fact-ub}, in the formulation
in~\eqref{eq:fact-ub}, shows that, for any $m\times n$ matrix $\mat$,
\(
\disc(\mat) \lesssim \gamma_2(\mat)\sqrt{\log(2m)},
\)
where
\begin{equation}\label{eq:ch7-gamma2-defn}
 \gamma_2(\mat) \eqdef \inf\{\|U\|_{2\to\infty}\|V\|_{1\to2}: UV = \mat\}.
\end{equation}
The notation $\gamma_2$ comes from functional analysis,
where this function was first studied in connection with factoring
linear operators between Banach spaces through a Hilbert space. 

The following corollary follows from the above and the observation
that $\gamma_2(\mat)$ is a hereditary upper bound on discrepancy.

\begin{corollary}[of Lemma~\ref{lm:fact-ub}]\label{cor:ch7-fact-ub} 
  For any $m\times n$ matrix $\mat$, we have
  \[
    \herdisc(\mat) \lesssim \sqrt{\ln(2m)} \gamma_2(\mat).
  \]
\end{corollary}
\begin{proof}
  This follows immediately from Lemma~\ref{lm:fact-ub}
  and the observation that $\gamma_2(\mat)$ is monotonically
  non-increasing under taking submatrices. Indeed, for any $J
  \subseteq [n]$, and any  factorization $\mat = UV$, we have that
  $\mat(*,J) = UV(*,J)$, and $\|V(*,J)\|_{1\to 2} \le \|V\|_{1\to 2}$. Minimizing over all choices of $U$ and $V$
  gives us $\gamma_2(\mat(*,J)) \le \gamma_2(\mat)$. Thus, by
  Lemma~\ref{lm:fact-ub}, we have
  \[
    \disc(\mat(*,J)) \lesssim \sqrt{\ln(2m)} \gamma_2(\mat(*,J))
    \le
    \sqrt{\ln(2m)} \gamma_2(\mat).
  \]
  Taking the maximum on the left hand side over all choices of $J\subseteq[n]$ gives the
  corollary. 
\end{proof}

Note that the all-ones matrix is no longer a bad example for
Corollary~\ref{cor:ch7-fact-ub}, since both the hereditary discrepancy
and the $\gamma_2$ norm of this matrix are equal to $1$. 

\begin{exercise}\label{ex:ch7-fact-ub}
  Give an example in which Corollary~\ref{cor:ch7-fact-ub} is tight up
  to constant factors. That is, show that for infinitely many values of
  $m$, there exists an $m\times n$ matrix $\mat$ for which
  $\disc(\mat) \gtrsim \sqrt{\ln(2m)} \gamma_2(\mat)$.
\end{exercise}

We now show that $\gamma_2(\mat)$ also gives a nearly tight
lower bound on hereditary discrepancy (see Theorem
\ref{thm:ch7-fact-lb} for the formal statement). This will be our goal in the rest of this section.

\subsection{Semidefinite Program Formulation}
%Before showing this, 
%in Corollary~\ref{cor:ch7-fact-ub}
%is tight up to a poly-logarithmic factor, 
We start  by reformulating
$\gamma_2(\mat)$ as a semidefinite program (SDP). One important
consequence of this formulation is that $\gamma_2(\mat)$ is
efficiently computable, and, therefore, so is the upper bound we
proved. By semidefinite programming duality, this formulation also gives a dual formulation of
$\gamma_2(\mat)$ as a maximization problem. This dual formulation will
be essential in giving a lower bound on $\herdisc(\mat)$ in terms of
$\gamma_2(\mat)$.

Recall the definition of $\gamma_2(\mat)$ in
\eqref{eq:ch7-gamma2-defn}. Note first that we can assume that the
infimum is over factorizations $\mat = UV$ for which
$\|U\|_{2\to\infty} = \|V\|_{1\to 2} = \sqrt{\gamma_2(\mat)}$. This is because, for any $c >
0$, we can replace $U$ by $cU$ and $V$ by $\frac1c V$ to get another
factorization, and choosing
$c \eqdef \sqrt{\frac{\|V\|_{1\to   2}}{\|U\|_{2\to\infty}}}$ makes
the two matrix norms equal.

Note also that, since $\gamma_2(\mat) \le \|I\|_{2\to \infty} \|\mat\|_{1\to 2}=  \|\mat\|_{1\to 2}$, we can
further assume that $\|U\|_{2\to\infty}$ and $\|V\|_{1\to 2}$ are both
bounded by $\|\mat\|_{1\to 2}^{1/2}$. Then it follows by a standard
compactness argument that the infimum in the definition of
$\gamma_2(\mat)$ is achieved. 

Let us now consider some factorization $\mat = UV$ of the $m\times n$
matrix $\mat$ into a $m \times k$ matrix $U$ and a $k\times n$ matrix
$V$. The matrix
\[
  X= \begin{pmatrix}
    U\\
    V^T
  \end{pmatrix}
  \begin{pmatrix}
    U^T  & V 
  \end{pmatrix}
  =
  \begin{pmatrix}
    UU^T & UV\\
    V^TU^T & V^T V
  \end{pmatrix}
  =
  \begin{pmatrix}
    UU^T & \mat\\
    \mat^T & V^T V
  \end{pmatrix}
\]
is positive semidefinite, since it equals $W^T W$ for
$W = \begin{pmatrix}U^T&  V\end{pmatrix}$. Notice that the entry $X(i,i)$, for 
$1 \le i \le m$, equals $\|U(i,*)\|_2^2$, and the entry $X(m + j,m+j)$  equals $\|V(*,j)\|_2^2$ for
$1 \le j \le n$.

Conversely, suppose that $X$
is a semidefinite $(m+n) \times (m+n)$ matrix such that the submatrix
formed by the first $m$ rows and last $n$ columns equals $\mat$. Since
$X$ is positive semidefinite, we can write it as $W^T W$ for some
$(m+n) \times (m+n)$ matrix $W$. Defining $U$ to be the transpose of
the submatrix of $W$ formed by taking its first $m$ columns, and $V$
to be the submatrix of $W$ formed by its last $n$ columns, we get a
factorization $\mat = UV$. Moreover,
$\max_{i = 1}^m X(i,i) = \|U\|_{2\to\infty}^2$, and
$\max_{j=1}^n X(m+j,m+j) = \|V\|_{1\to 2}^2$. %With these observations
%in hand, we can prove the following equivalence.
This equivalence directly gives the following lemma.
%detailed proof we leave as an exercise. 
\begin{lemma}\label{lm:ch7-fact-sdp}
  For any $m\times n$ matrix $\mat$ we have
  \begin{align}
    \gamma_2(\mat) = &\min t \label{eq:ch7-sdp-obj} \\
               %      &\text{subject to}\notag\\
                 \text{s.t. \,\,\, }    &X(i,m+j) = X(m+j,i) =\mat(i,j) \ \ \forall i\in [m], j \in [n]
                    \label{eq:ch7-sdp-eq}\\
                     &X(i,i) \le t\ \ \forall i \in [m+n]\\
                     &X \succeq 0\label{eq:ch7-sdp-psd}
  \end{align}
  where $X$ ranges over symmetric $(m+n) \times (m+n)$
  real matrices. 
\end{lemma}
% \begin{proof}
%   The observations above already give us that
%   \begin{align}
%     \gamma_2(\mat) = &\inf \left( \max_{i = 1}^m X(i,i)\right)^{1/2}\left(\max_{j=1}^n X(m+j,m+j)\right)^{1/2}\label{eq:ch7-nl-obj}\\
%                      &\text{subject to}\notag\\
%                      &X(i,m+j) = X(m+j,i) =\mat(i,j) \ \ \forall i\in [m], j \in [n]\\
%                      &X \succeq 0\label{eq:ch7-nl-psd}
%   \end{align}
%   Since the arithmetic mean of two non-negative numbers is at most
%   their maximum,
%   \begin{multline}\label{eq:ch7-norms-equal}
%     \left( \max_{i = 1}^m X(i,i)\right)^{1/2}\left(\max_{j=1}^n X(m+j,m+j)\right)^{1/2}\\
%     \le
%     \max\left\{ \max_{i = 1}^m X(i,i),\max_{j=1}^n X(m+j,m+j)\right\}
%     =   \max_{i = 1}^{m+n} X(i,i)
%   \end{multline}
%   We can show that this inequality can be assumed to hold with
%   equality. Let $c > 0$ be chosen later, and let $D$ be a diagonal matrix such that $D(i,i) = c$ for $i
%   \in [m]$, and $D(m+j,m+j) = \frac1c$ for $j \in [n]$. Then if $X$ is
%   positive semidefinite, $DXD$ is as well. Moreover, notice that, for
%   any $i \in [m]$ and $j \in [n]$
%   \[
%     (DXD)(i,m+j) = (DXD)(m+j,i) = c\cdot\frac1c X(i,m+j) = X(i,m+j).
%   \]
%   Therefore, for any feasible solution $X$ to
%   \eqref{eq:ch7-nl-obj}--\eqref{eq:ch7-nl-psd}, $DXD$ is also
%   feasible. Furtheremore,
%   \(
%   \max_{i = 1}^m (DXD)(i,i)  =   c^2\max_{i = 1}^m X(i,i),
%   \)
%   and
%   \(
%   \max_{j = 1}^n (DXD)(m+j,m+j)  =   \frac{1}{c^2}\max_{j = 1}^n X(m+j,m+j),
%   \)
%   so $DXD$ also achieves the same objective value in
%   \eqref{eq:ch7-nl-obj} as $X$. 
%   Choosing $c = \sqrt{\frac{\max_{j  = 1}^n X(m+j,m+j)}{\max_{i = 1}^m
%       X(i,i)}}$ and replacing $X$ by $DXD$ now allows us to assume that $\max_{i = 1}^m
%   X(i,i)= \max_{j=1}^n X(m+j,m+j)$, in which case inequality \eqref{eq:ch7-norms-equal}
%   holds with equality. Therefore we get
%   \begin{align*}
%     \gamma_2(\mat) = &\inf \max_{i = 1}^{m+n} X(i,i)\\
%                      &\text{subject to}\notag\\
%                      &X(i,m+j) = X(m+j,i) = \mat(i,j) \ \ \forall i\in [m], j \in [n]\\
%                      &X \succeq 0
%   \end{align*}
%   This is easily seen to be equivalent to
%   \eqref{eq:ch7-sdp-obj}--\eqref{eq:ch7-sdp-psd} except for taking an
%   $\inf$ rather than a $\min$.% : once $X$ is fixed, the optimal choice
%   % for the $t$ variable in
%   % \eqref{eq:ch7-sdp-obj}--\eqref{eq:ch7-sdp-psd} is the maximum of
%   % $\max_{i = 1}^m X(i,i)$ and $\max_{j=1}^n X(m+j,m+j)$.
%   To see that
%   the infimum is also achieved, notice that
%   $\gamma_2(\mat) \le \|\mat\|_{1\to 2}$ because we can always take
%   the trivial factorization $\mat = I\mat$, with $I$ the $m \times m$
%   identity matrix. This means we only need to take the infimum over
%   $X \succeq 0$ satisfying $\max_{i=1}^m X(i,i) \le \|\mat\|_{1\to 2}$
%   as well as the equalities \eqref{eq:ch7-sdp-eq}. Since this set is
%   compact, the infimum is achieved.
% \end{proof}


Semidefinite programs share two important properties with linear
programs: they are efficiently solvable, in the sense that an
approximately optimal solution can be found in polynomial time to any
degree of accuracy; they satisfy a duality theory, i.e., we can 
derive a dual maximization semidefinite program to a given
minimization semidefinite program, so that the optimal values of the
two programs are equal.\footnote{These statements are imprecise: in
  order to be efficiently solvable and satisfy strong duality,
  semidefinite programs need to be sufficiently ``nice''. This turns
  out not to be an issue for our application.} Then, using
Lemma~\ref{lm:ch7-fact-sdp}, we can show that
$\gamma_2(\mat)$ is efficiently computable, and has a dual
maximization formulation. These facts, whose proofs we omit, are
captured in the following lemma.

%\sasho{might add a proof and more about duality of SDPs here}

\begin{lemma}\label{lm:ch7-fact-dual}
  For any $m\times n$ matrix $\mat$, and any $\varepsilon > 0$, we can
  compute a factorization $\mat = UV$ so that $\|U\|_{2\to\infty}
  \|V\|_{1\to 2} \le (1+\varepsilon)\gamma_2(\mat)$ in time polynomial
  in $m$, $n$, the maximum bitsize of any entry of $\mat$, and in
  $\log(1/\varepsilon)$.

  Moreover, $\gamma_2(\mat)$ satisfies the dual formula
  \begin{equation}
    \label{eq:ch7-fact-dual}
    \gamma_2(\mat) =
    \max\{\|P\mat Q\|_{\tr}: P, Q \text{ diagonal }, \tr(P^2) = \tr(Q^2) = 1\}.
  \end{equation}
  Above, $\|P\mat Q\|_{\tr}$ is the trace norm of the matrix $P\mat Q$,
  i.e., the sum of its singular values. 
\end{lemma}

\subsection{Lower Bound}
%We will now show that the upper bound in
%Corollary~\ref{cor:ch7-fact-ub} is tight up to logarithmic factors. 
To lower bound the hereditary discrepancy by $\gamma_2(\mat)$ up to logarithmic factors, we need a general tool for proving lower bounds on hereditary
discrepancy. 
We next define an intermediate quantity called the determinant lower bound $\detlb(\mat)$. We use the connection between rounding and hereditary discrepancy, discussed Chapter~\ref{ch1-intro}, and a covering argument, to show that $\detlb(\mat)$ is a lower bound on hereditary discrepancy. Finally, we relate  $\gamma_2(\mat)$ and $\detlb(\mat)$.

\paragraph{The Determinant Lower bound and Hereditary Discrepancy.}
The determinant lower bound
  of an $m\times n$ matrix $\mat$ is defined as
  \begin{equation}
    \label{eq:ch7-detlb-defn}
    \detlb(\mat) \eqdef \max_k \max_{B \in S_k(\mat)} |\det(B)|^{1/k},
  \end{equation}
  where $S_k(\mat)$ is the set of $k\times k$ submatrices of
  $\mat$.
We have the following lower bound on $\herdisc(\mat)$.
\begin{lemma}\label{lm:ch7-detlb}
  For any $m\times n$ matrix $\mat$ we have the inequality
  \[
    \herdisc(\mat) \ge \frac12 \detlb(\mat).
  \]
\end{lemma}
\begin{proof}
We prove a slightly weaker bound
$\herdisc(\mat) \ge  \detlb(\mat)/4$,
  and leave the tighter bound as an exercise.
  We show that for any $k\times k$ square matrix $B$,
  \(
  \herdisc(B) \ge (1/4) |\det(B)|^{1/k}.
  \)
  The claimed bound then follows directly as $\herdisc(\mat) \ge
  \max_k \max_{B \in S_k(\mat)} \herdisc(B)$.   
  
\smallskip
  It follows from the transference theorem (Theorem~\ref{thm:ch1-transference}) that 
  \[
    \forall y \in [0,1]^k \ \exists z \in \{0,1\}^k\ \ 
  \|B (y-z)\|_\infty \le \herdisc(B).
  \]
  Let $K \eqdef \{x \in \R^k: \|Bx\|_{\infty} \le
  \herdisc(B)\}$. The equation above is equivalent to
  \(
    [0,1]^k \subseteq \bigcup_{z \in \{0,1\}^k} (K + z).
  \)
  Letting $\vol^k(\cdot)$ be the $k$-dimensional Lebesgue
  measure, by its translation invariance we have
  \begin{align}
    1 &= \vol^k([0,1]^k) \le
    \sum_{z \in \{0,1\}^k} \vol^k(K + z) = 2^k \vol^k(K).\label{eq:ch7-covering}
  \end{align}
  To turn this into a lower bound on $\herdisc(B)$, we need to
  relate it to $\vol^k(K)$. Let us assume that $B$ is invertible: 
  otherwise, the lower bound
  \(
  \herdisc(B) \ge \frac14 |\det(B)|^{1/k}=0
  \)
  is trivial. Then we can write
  \begin{align*}
    K & = \{x \in \R^k: \|B x\|_{\infty} \le   \herdisc(B)\}\\
    & = \{B^{-1} y: y\in \R^k, \|y\|_{\infty} \le \herdisc(B)\}.
  \end{align*}
  The right side is the image of the cube
  $[-\herdisc(B),\herdisc(B)]^k$ under multiplication by $B^{-1}$, and
  has volume $
  (2\herdisc(B))^{k} |\det(B^{-1})|$ $=
  2^k\herdisc(B)^k/|\det(B)|$.
  
Together with \eqref{eq:ch7-covering}, 
 we get that \[
    1 \le 2^k\vol^k(K) = 4^k\herdisc(B)^k/|\det(B)|. \]
  Re-arranging gives us the lower bound. 
\end{proof}

\begin{remark}\label{rem:vollb}
  The determinant lower bound can be far from tight when $m$ is
  large. For example, take $\mat$ to be a $2^n \times n$ matrix whose
  rows are all the vectors in $\{-1, +1\}^n$. Hadamard's inequality shows that
  $\detlb(\mat) \le \sqrt{n}$, but it is easy to see that $\disc(\mat)
  = n$. Intuitively, the reason for this gap is that the determinant
  lower bound only considers square submatrices of $\mat$. The proof
  above, however, can be modified slightly to establish a
  stronger volumetric lower bound on hereditary discrepancy, as we
  explain next. 

  Notice that the proof of the inequality \eqref{eq:ch7-covering} did
  not use that $B$ is a square matrix, and in fact works for any
  submatrix $B$ of $\mat$. Using the inequality for
  submatrices that consist of a subset of the columns of $\mat$, we
  get that, for any $S \subseteq [n]$, the convex body
  $\{x \in \R^S: \|\mat(*,S) x\|_{\infty} \le \herdisc(\mat)\}$ has
  $|S|$-dimensional volume at least $2^{-|S|}$, or, equivalently,
  $\{x \in \R^S: \|\mat(*,S) x\|_{\infty} \le 1\}$, has volume at
  least $(2\herdisc(\mat))^{-|S|}$. Letting 
  $\overline{K} \eqdef \{x \in \R^n: \|\mat x\|_{\infty} \le 1\}$,
  this means that
  \begin{equation}
    \label{eq:vollb}
    \herdisc(\mat) \ge \frac12 \vollb(\mat) \eqdef
    \min_{S \subseteq [n]}\vol^{S}(\overline{K} \cap \R^S)^{-1/|S|},
  \end{equation}
  where $\vol^S(\cdot)$ is the $|S|$-dimensional Lebesgue measure
  taken inside $\R^S$, and we used the observation that
  \(
  \overline{K} \cap \R^S = \{x \in \R^S: \|\mat(*,S) x\|_{\infty} \le 1\}.
  \)
   This volumetric lower bound takes advantage of
  ``tall'' matrices, and can be stronger than the determinant lower
  bound. Returning to the $2^n \times n$ matrix $\mat$ above, whose rows are
  all vectors in $\{-1,+1\}^n$, we have
  $\overline{K} = \{x \in \R^n: \|x\|_1 \le 1\}$, i.e., the
  $n$-dimensional $\ell_1$ ball, which has Lebesgue measure
  $2^n/n!$. So $\vollb(\mat) = (n!)^{1/n}/2 \approx n/2e$ by
  Stirling's approximation. Unlike the determinant lower bound, this
  matches the discrepancy of $\mat$  up to a constant factor.
\end{remark}

\begin{exercise}
  Prove the inequality \(\herdisc(\mat) \ge (1/2) \detlb(\mat)\) in
  Lemma~\ref{lm:ch7-detlb}. To do so, argue that
  \eqref{eq:ch7-covering} can be strengthened to $\vol^k(K) \ge 1$.
\end{exercise}

\paragraph{Relating $\detlb(\mat)$ and $\gamma_2(\mat)$.}
%We show a lower bound on $\herdisc(\mat)$ in terms of $\gamma_2(\mat)$
Our goal will be to show that $\gamma_2(\mat)\lesssim \log(2n) \detlb(\mat)$. 

Let us first see how to use
Lemma~\ref{lm:ch7-fact-dual} to get an inequality in the other
direction. Let us take any $k\times k$ submatrix $\mat(S,T)$ of $\mat$ that
achieves $\detlb(\mat)$. Define $P(i,i) = \frac{1}{\sqrt{k}}$ for $i \in S$,
$Q(j,j) = \frac{1}{\sqrt{k}}$ for $j \in T$, and define all other entries of $P$
and $Q$ to be $0$. Then, by \eqref{eq:ch7-fact-dual},
\[
  \gamma_2(\mat) \ge \|P\mat Q\|_{\tr} = \frac1k \|\mat(S,T)\|_{\tr},
\]
and notice that $\frac1k \|\mat(S,T)\|_{\tr}$ is the arithmetic mean of
the singular values of $\mat(S,T)$. Since $|\det(\mat(S,T))|^{1/k}$ is
the geometric mean of the singular values of the same matrix, the
AM-GM inequality gives us that
$\gamma_2(\mat) \ge \detlb(\mat)$.  

To reverse this inequality, we
then need an inequality that reverses the AM-GM inequality, in the
appropriate sense. The next lemma is sufficient for us for this
purpose. %\sasho{Is this ``reverses AM-GM'' confusing?}

\begin{lemma}\label{lm:ch7-rev-AMGM}
  %For any integer $r \ge 2$, and 
  For any non-negative real numbers $\sigma_1 \ge \ldots \ge
  \sigma_r \ge 0$,
  \[
    \max_{k = 1}^r k(\sigma_1 \cdots \sigma_k)^{1/k}
    \ge
    \max_{k = 1}^r k\sigma_k
    \ge
    \frac{\sigma_1 + \ldots + \sigma_r}{1+\ln r}
  \]
\end{lemma}
\begin{proof}
  The first inequality is trivial. The second
  follows as,
  %By the (the trivial case of) H\"older's inequality, we get that 
  \begin{align*}
    \sum_{k = 1}^r \sigma_k =
    \sum_{k = 1}^r k\sigma_k\cdot\frac{1}{k}
    &\le \left(\max_{k = 1}^r k\sigma_k\right)\left(\sum_{k=1}^r \frac{1}{k}\right)
    \le \left(\max_{k = 1}^r k\sigma_k\right)(1+\ln r),
  \end{align*}
  %where we used the inequality
  as $\sum_{k=1}^r \frac{1}{k} \le 1 + \ln r$.
\end{proof}

Let us also recall the Cauchy-Binet formula, which states that for any
two $m\times n$ matrices $A,B$, where $n \le m$, we have
\[
  \det(A^T B) = \sum_{S \subseteq [m]: |S| = n} \det(A(S,*))\det(B(S,*)).
\]
For a short proof, see Section 3.2 in Tao's
book~\cite{Tao-randmat}. Note that applying the Cauchy-Schwarz
inequality to the right hand side of the Cauchy-Binet formula, we get
\begin{align}
  \det(A^T B)^2 &\le
  %\left(\sum_{S \subseteq [m]: |S| = n} \det(A(S,*))\det(A(S,*)\right)\left(\sum_{S \subseteq [m]: |S| = n} \det(B(S,*))\det(B(S,*)\right)\notag\\
  %&=
     \det(A^T A) \det(B^T B).  \label{eq:ch7-CS-mat}
\end{align}

We need one final lemma before we can relate $\gamma_2(\mat)$ and
$\detlb(\mat)$. This lemma  allows us to extract a submatrix of
$\mat$ to use as a certificate that $\detlb(\mat)$ is large. 

\begin{lemma}\label{lm:ch7-submatrix}
  Let $n\le m$ be positive integers, and let $D$ be an $m\times m$ diagonal matrix with non-negative entries such that $\tr(D) =
  1$. Then for any $m\times n$ matrix $A$, there exists a set
  $S\subseteq [m]$ of cardinality $n$ for which 
  \[
    |\det(A(S,*))| \ge  \sqrt{n!\det(A^TDA)}.
  \]
\end{lemma}
\begin{proof}
  By the Cauchy-Binet formula and H\"older's inequality, we have
  \begin{align*}
    \det(A^T DA) %&= \sum_{S\subseteq [m]: |S| = n} \det(D(S)\det(A(S,*))^2\\
                 &= \sum_{S\subseteq [m]: |S| = n}  \det(A(S,*))^2 \left(\prod_{i \in S} D(i,i)\right)\\
                 &\le \left(\max_{S \subseteq [m]: |S| = n}\det(A(S,*))^2\right)
                   \left(\sum_{S\subseteq [m]: |S| = n} \prod_{i \in S} D(i,i)\right).
  \end{align*}
  It remains to bound the sum on the right hand side. One can show that this sum is maximized when all $D(i,i)$ are equal,
  in which case it equals ${m \choose n} \frac{1}{m^n} <
  \frac{1}{n!}$. For an elementary proof of a similar bound,
  notice that, since all $D(i,i)$ are non-negative,
  \[
    1 = (D(1,1) + \ldots + D(m,m))^n
    \ge n! \sum_{S\subseteq [m]: |S| = n} \prod_{i \in S} D(i,i),
  \]
  because when we expand $(D(1,1) + \ldots + D(m,m))^n$ we get each
  multilinear term $\left(\prod_{i \in S} D(i,i)\right)$ $n!$ times,
  in addition to the other terms. %Combining the inequalities gives the lemma.
\end{proof}

We are, finally, ready to prove that $\gamma_2(\mat)$ and
$\detlb(\mat)$ are equal up to a logarithmic factor.

\begin{lemma}\label{lm:ch7-fact-detlb}
  For any $m\times n$ matrix $\mat$ of rank $r$, we have
  \[
    \detlb(\mat) \le \gamma_2(\mat) \lesssim (1 + \log r) \detlb(\mat).
  \]
\end{lemma}
\begin{proof}
  As we argued above, the first inequality follows by considering the
  $k\times k$ submatrix $\mat(S,T)$ of $\mat$ achieving
  $|\det(\mat(S,T))|^{1/k} = \detlb(\mat)$, and setting $P$ and $Q$ in
  \eqref{eq:ch7-fact-dual} so that $P\mat Q = \frac1k \mat(S,T)$.
  Here we prove the second inequality, again using
  Lemma~\ref{lm:ch7-fact-dual}, as well as
  Lemmas~\ref{lm:ch7-rev-AMGM}~and~\ref{lm:ch7-submatrix}.

  Let $P$ and $Q$ achieve equality in \eqref{eq:ch7-fact-dual}, and
  let $\sigma_1 \ge \ldots \ge \sigma_r \ge 0$ be the $r$ largest singular
  values of $P \mat Q$. Since $P \mat Q$ has rank at most $r$, all
  other singular values of $P \mat Q$ are $0$, and $\gamma_2(\mat) =
  \sigma_1 + \ldots + \sigma_r$. Let $k \in [r]$ be the integer for
  which Lemma~\ref{lm:ch7-rev-AMGM} guarantees that
  \begin{equation}\label{eq:ch7-top-k}
    k(\sigma_1 \cdots \sigma_k)^{1/k}
    \ge \frac{\sigma_1 + \ldots + \sigma_r}{1+\ln r}
    = \frac{\gamma_2(\mat)}{1+\ln r}.
  \end{equation}
  The goal in the rest of the proof is to show that there exists a
  $k\times k$ submatrix $\mat(S,T)$ of $\mat$ for which $|\det(\mat(S,T))|^{1/k}$ is not
  much smaller than $k(\sigma_1 \cdots \sigma_k)^{1/k}$.

  Let $v_1, \ldots, v_k$ be the right singular vectors of $P \mat Q$
  associated with $\sigma_1, \ldots, \sigma_k$, and let $V$ be a
  matrix with columns $v_1, \ldots, v_k$. Define the $m\times k$
  matrix $A \eqdef \mat QV$. The matrix $PA$ then has singular values
  $\sigma_1, \ldots, \sigma_k$, or, equivalently, $A^TP^2A$ has
  eigenvalues $\sigma_1^2, \ldots,
  \sigma_k^2$. Lemma~\ref{lm:ch7-submatrix} applied to $A$ and $D
  \eqdef P^2$ now tells us that there exists a set $S \subseteq [m]$
  of cardinality $k$ so that
  \begin{align}
    |\det(\mat(S,*)QV)| &= |\det(A(S,*))|\notag\\
    &\ge \sqrt{k! \det(A^TP^2A)} = \sqrt{k!} \sigma_1 \cdots \sigma_k.    \label{eq:ch7-rows}
  \end{align}

  Next we would like to apply Lemma~\ref{lm:ch7-submatrix} again, in
  order to replace $\mat(S,*)QV$ with $\mat(S,T)$ for some cardinality
  $k$ set $T \subseteq [n]$. We first need to deal with the matrix
  $V$, however. To this end, apply \eqref{eq:ch7-CS-mat}  to the
  matrices $Q\mat(S,*)^T$ and $V$, to conclude that
  \begin{align}
    \det(\mat(S,*)Q^2 \mat(S,*)^T) &\ge
    \frac{\det(\mat(S,*)QV)^2}{\det(V^T V)} \notag\\
                                   &= \det(\mat(S,*)QV)^2,\label{eq:ch7-proj}
  \end{align}
  where we used that $V^T V = I$. 
  Then 
  \eqref{eq:ch7-rows} and \eqref{eq:ch7-proj} together give us
  \[
    \sqrt{\det(\mat(S,*)Q^2\mat(S,*)^T)} \ge
    \sqrt{k!} \sigma_1 \cdots \sigma_k.
  \]
  Finally, we apply Lemma~\ref{lm:ch7-submatrix} one more time, now
  with the $n\times k$ matrix $A \eqdef \mat(S,*)^T$ and $D \eqdef
  Q^2$, concluding that there exists a set $T \subseteq [n]$ of
  cardinality $k$ for which
  \[
    |\mat(S,T)| \ge \sqrt{k!\det(\mat(S,*)Q^2\mat(S,*)^T)} \ge
    k! \sigma_1 \cdots \sigma_k.
  \]
  Together with \eqref{eq:ch7-top-k}, and as $k! \ge (k/e)^k$, we can conclude that
  \begin{align*}
    \detlb(\mat)
    \ge
    |\det(\mat(S,T))|^{1/k}
   % &\ge
    %  (k! \sigma_1 \cdots \sigma_k)^{1/k}\\
    &\gtrsim
      k (\sigma_1 \cdots \sigma_k)^{1/k}
    \ge
      \frac{\gamma_2(\mat)}{1+\ln r}.  \qedhere
  \end{align*}
%  as we needed to prove. Above, on the second line, we used the
%  inequality 
\end{proof}

Lemma~\ref{lm:ch7-fact-detlb} gives an approximation algorithm for the
determinant lower bound: computing $\gamma_2(\mat)$ gives a factor
$O(\log r)$ approximation for $\detlb(\mat)$ for any rank $r$ matrix
$\mat$. No better polynomial time computable approximation for
$\detlb(\mat)$ is known. 

Combining Lemmas~\ref{lm:ch7-detlb}~and~\ref{lm:ch7-fact-detlb}, we get
the following lower bound on hereditary discrepancy.

\begin{theorem}\label{thm:ch7-fact-lb}
  For any $m\times n$ matrix $\mat$ of rank $r$, we have
  \[
    \herdisc(\mat) \gtrsim  \frac{\gamma_2(\mat)}{1 + \log r}.
  \]
\end{theorem}


\subsection{The Approximation Theorem}

We can now state our main result.
\begin{theorem}\label{thm:ch7-fact-main}
  For any $m\times n$ matrix $\mat$ of rank $r$, we have
  \[
    \frac{\gamma_2(\mat)}{1 + \log r} \lesssim
    \herdisc(\mat)
    \lesssim
    \sqrt{\ln(2m)} \gamma_2(\mat).
  \]
\end{theorem}
Theorem~\ref{thm:ch7-fact-main} follows immediately from
Corollary~\ref{cor:ch7-fact-ub}, and Theorem~\ref{thm:ch7-fact-lb}.

In Exercise~\ref{ex:ch7-fact-ub} we showed that the upper bound in
Theorem~\ref{thm:ch7-fact-main} is tight up to constant factors. Soon
we will see a simple example of a matrix for which the lower bound is
tight.

\section{Some Properties}

There are many natural questions one can ask about the behavior of
hereditary discrepancy under simple transformations. For example, we
can ask how $\herdisc(\SS \cup \SS')$ compares to $\herdisc(\SS)$ and
$\herdisc(\SS')$ for two set systems $\SS$ and $\SS'$ on the same
universe. As another example, we can ask how $\herdisc(\mat^T)$
compares to $\herdisc(\mat)$ for a matrix $\mat$. None of these
questions has an obvious answer, but, nevertheless, we can give nearly tight
answers to them using Theorem~\ref{thm:ch7-fact-main}, and simple
properties of the $\gamma_2$ norm. Let us remark that it is essential
here to work with hereditary discrepancy - discrepancy itself does not
behave nicely with respect to most simple transformations. Two
examples of this are offered in the next exercises.

\begin{exercise}
  Give an example of two set systems $\SS$ and $\SS'$, both on the ground
  set $[n]$ for an arbitrarily large integer $n$, so that $\disc(\SS)
  = \disc(\SS') = 0$ but $\disc(\SS \cup \SS') \gtrsim n$.
\end{exercise}

% SS: each set takes t elements from L := {1, ..., n/2} and t elements from
% R := {n/2 + 1, ..., n}. SS contains all such sets.
% SS' = {L, R}.
% e.g. t= n/4 works

\begin{exercise}
  Give an example of a $n\times n$ matrix with $\{0,1\}$ entries,
  for $n$ arbitrarily large, so that $\disc(\mat) = 0$ but
  $\disc(\mat^T) \gtrsim \sqrt{n}$.
\end{exercise}

% Take a Hadamard matrix, replace -1 with 0, duplicate each columns
% twice, pad with 0 rows to make it square.

The next lemma summarizes the basic properties of the $\gamma_2$ norm
that are useful in reasoning about hereditary
discrepancy. Below, we use $\otimes$ for the Kronecker product: given
two matrices $A$ and $B$, of dimensions, respectively, $m \times n$
and $p \times q$, $A\otimes B$ is a $(mp) \times (nq)$ matrix whose
rows are indexed by $[m]\times [p]$, whose columns are indexed by $[n]\times
[q]$, and whose entries are $(A\otimes B)((i_1, i_2), (j_1, j_2)) =
A(i_1,j_1)B(i_2,j_2)$. The key property of Kronecker products that we
repeatedly use is that $(A\otimes B)(A'\otimes B') =
(AA')\otimes(BB')$ for any matrices $A,A',B,B'$ for which these matrix
products are well-defined.

\begin{lemma}\label{lm:ch7-gamma2-props}
  The $\gamma_2$ function on matrices has the following properties:
  \begin{enumerate}
  \item $\gamma_2(\mat(S,T)) \le \gamma_2(\mat)$ for any $m\times n$
    matrix $\mat$ and any submatrix $\mat(S,T)$ of
    it;\label{pr:ch7-submatrix}
  \item duplicating rows and columns of $\mat$  keeps $\gamma_2(\mat)$ unchanged;\label{pr:ch7-duplicate}
  \item $\gamma_2(\mat) = \gamma_2(\mat^T)$ for any matrix $\mat$;\label{pr:ch7-transpose}
  \item $\gamma_2(\mat) = 0 \implies \mat = 0$;\label{pr:ch7-pos}
  \item $\gamma_2(c\mat) = |c|\gamma_2(\mat)$ for any matrix $\mat$
    and real value $c$;\label{pr:ch7-hom}
  \item $\gamma_2(\mat + \mat') \le \gamma_2(\mat) + \gamma_2(\mat')$;\label{pr:ch7-triangle}
  \item for any two matrices $\mat'$ and $\mat''$, each with $n$ columns, and the matrix 
    $\mat = \begin{pmatrix}\mat'\\ \mat''\end{pmatrix}$, 
    we have $\gamma_2(\mat) \le \sqrt{\gamma_2(\mat')^2 +
      \gamma_2(\mat'')^2}$; \label{pr:ch7-union}
  \item $\gamma_2(\mat \otimes \mat') = \gamma_2(\mat)
    \gamma_2(\mat')$ for any matrices $\mat$ and $\mat'$.\label{pr:ch7-tensor}
  \end{enumerate}
\end{lemma}
\begin{proof}
  Property~\ref{pr:ch7-submatrix} was noted previously. Next we
  give the proofs of some of the remaining properties, and leave
  the rest as exercises.

  Property~\ref{pr:ch7-triangle} can be proved in several ways,
  including by directly giving a factorization of $\mat + \mat'$ using
  optimal factorizations of $\mat$ and $\mat'$. Instead, let us prove
  it using the SDP characterization from
  Lemma~\ref{lm:ch7-fact-sdp}. Let $X, t$ be a solution to the
  semidefinite program \eqref{eq:ch7-sdp-obj}--\eqref{eq:ch7-sdp-psd}
  for $\mat$ that achieves the optimal objective value
  $\gamma_2(\mat)$, and let $X', t'$ be a solution for $\mat'$ that
  achieves objective value $\gamma_2(\mat')$. Then $X + X', t + t'$ is a
  feasible solution for $\mat + \mat'$, and achieves objective value
  $\gamma_2(\mat) + \gamma_2(\mat')$, which is an upper bound $\gamma_2(\mat + \mat')$.

\smallskip
  While property~\ref{pr:ch7-union} can also be shown using the SDP
  characterization, we give a direct
  factorization of $\mat = \begin{pmatrix}\mat'\\ \mat''\end{pmatrix}$
  instead. Let $\mat' = U' V'$ and $\mat'' = U'' V''$ be
  factorizations achieving, respectively,
  $\gamma_2(\mat') = \|U'\|_{2\to\infty} \|V'\|_{1\to 2}$, and
  $\gamma_2(\mat'') = \|U''\|_{2\to\infty} \|V''\|_{1\to
    2}$. 
    %Moreover, choose $U'$ and $U''$ 
    Rescale, so that
  $\|U'\|_{2\to\infty} = \|U''\|_{2\to\infty} = 1$,
  $\|V'\|_{1\to 2} = \gamma_2(\mat')$, and
  $\|V''\|_{1\to 2} = \gamma_2(\mat'')$.
  %This is always possible to
  %achieve by replacing $U'$ by $\frac{U'}{\|U'\|_{2\to\infty}}$, and
  %$V'$ by $\|U'\|_{2\to\infty} V'$, and similarly for $U''$ and
  %$V''$. 
  Consider then the factorization $\mat = UV$ where
  \[
    U =
    \begin{pmatrix}
      U' &0\\
      0  &U''
    \end{pmatrix}; \ \ \ \ 
    V =
    \begin{pmatrix}
      V'\\
      V''
    \end{pmatrix}.
  \]
  %where the $0$ blocks in the definition of $U$ are of the appropriate
  %dimensions for $U$ to be well defined.
  Then, clearly,
  \(
  \|U\|_{2 \to \infty}= 1.
  \)
  %At the same time, since, by Pythagoras's theorem,
  Moreover, as
  \(\|V(*,i)\|_2^2 = \|V'(*,i)\|_2^2 + \|V''(*,i)\|_2^2\)
  for any $i \in [n]$, we have
  \begin{align*}
    \|V\|_{1\to 2}^2 &= \max_{i = 1}^n \left( \|V'(*,i)\|_2^2 + \|V''(*,i)\|_2^2 \right)\\
    &\le
      \max_{i = 1}^n  \|V'(*,i)\|_2^2 + \max_{i=1}^n\|V''(*,i)\|_2^2\\
    &=
    \|V'\|_{1\to 2}^2 + \|V''\|_{1\to 2}^2.
  \end{align*}
  Thus,
  \(
  \gamma_2(\mat)^2 \le \|U\|^2_{2\to\infty} \|V\|^2_{1\to 2}
  \le
  \gamma_2(\mat')^2 + \gamma_2(\mat'')^2,
  \)
  as claimed.
  
\smallskip
  Finally, we prove property \ref{pr:ch7-tensor}.  
  To
  show
  $\gamma_2(\mat \otimes \mat') \le \gamma_2(\mat) \gamma_2(\mat')$ we
  can, once again, either use the SDP characterization, or give a
  direct factorization. Using the SDP characterization and Lemma~\ref{lm:ch7-fact-sdp} let $X, t$ be
  a solution to the semidefinite program
  \eqref{eq:ch7-sdp-obj}--\eqref{eq:ch7-sdp-psd} for $\mat$ achieving
  value $\gamma_2(\mat)$, and let $X', t'$ be
  a solution for $\mat'$ achieving value $\gamma_2(\mat')$. Then $X\otimes X', tt'$ is a
  feasible solution for $\mat\otimes \mat'$ achieving value
  $\gamma_2(\mat)\gamma_2(\mat')$. Here we use the fact that the
  Kronecker product of two positive semidefinite matrices is positive
  semidefinite. 
  
  To show the reverse inequality $\gamma_2(\mat\otimes
  \mat') \ge \gamma_2(\mat)\gamma_2(\mat')$, we use
  Lemma~\ref{lm:ch7-fact-dual}. Let $P,Q$ achieve equality in
  \eqref{eq:ch7-fact-dual} for $\mat$, let $P', Q'$ achieve
  equality in \eqref{eq:ch7-fact-dual} for $\mat'$. Then,
  \begin{align*}
    \gamma_2(\mat\otimes \mat') &\ge
    \|(P\otimes P') (\mat \otimes \mat') (Q\otimes Q')\|_{\tr}\\
    &=
    \|(P\mat Q)\otimes (P'\mat Q')\|_{\tr}
    =
   \|P\mat Q\|_{\tr}\|P'\mat Q'\|_{\tr}.
 \end{align*}
 Above we used the fact that, for any matrices $A$ and $B$,
 $\|A\otimes B\|_{\tr} = \|A\|_{\tr}\|B\|_{\tr}$, which can be verified
 easily using the singular value decomposition. 
\end{proof}

\begin{exercise}
  Prove the remaining properties
  \ref{pr:ch7-transpose}--\ref{pr:ch7-hom} in Lemma~\ref{lm:ch7-gamma2-props}.
\end{exercise}

Note that properties \ref{pr:ch7-pos}--\ref{pr:ch7-triangle} in
Lemma~\ref{lm:ch7-gamma2-props} show that $\gamma_2$ is, indeed, a
norm on matrices.

Lemma~\ref{lm:ch7-gamma2-props} and Theorem~\ref{thm:ch7-fact-main}
show that properties~\ref{pr:ch7-submatrix}--\ref{pr:ch7-tensor} carry
over from the $\gamma_2$ norm to hereditary discrepancy, up to
logarithmic factors. Properties \ref{pr:ch7-submatrix},
\ref{pr:ch7-pos}, and \ref{pr:ch7-hom} hold trivially for
hereditary discrepancy even without the extra logarithmic factors, but
we know no direct proof for the others. The next theorem rephrases
some of these properties for set systems.

\begin{theorem}\label{thm:ch7-props-ss}
  The following bounds hold:
  \begin{enumerate}
  \item \label{pr:ch7-ss-dual} For a set system $(\SS, \uni)$, define the dual set system
    $\SS^*$ on the universe $\SS$ by $\{S_\elem: \elem \in \uni\}$ where $S_\elem = \{S \in
    \SS: \elem \in S\}$; then 
    \(
    \herdisc(\SS^*) \lesssim  \log(2|\uni|)^{3/2}\herdisc(\SS),
    \)
  \item \label{pr:ch7-ss-union} For any set systems $\SS_1, \ldots, \SS_k$ defined on the same
    universe $\uni$, 
 \[
 %\begin{multline*}
    \herdisc(\SS_1 \cup \cdots \cup \SS_k)
    \lesssim
    \log(2m)^{3/2} \left(\sum_{i=1}^k \herdisc(\SS_i)^2\right)^{1/2},\]
  %  \end{multline*}
    where $m = \sum_{i=1}^k|\SS_i|$.
    
  \item \label{pr:ch7-ss-product}  for set systems $(\SS_1, \uni_1), \ldots, (\SS_k, \uni_k)$,
    let $\SS_1 \times \cdots \times \SS_k$ consist of all product sets $S_1
    \times \cdots \times S_k$ with $S_i \in \SS_i$; then
    \begin{multline*}
      \log(2m)^{ - \frac{k}{2}-1}
      \lesssim
      \frac{\herdisc(\SS_1 \times \cdots \times \SS_k)}{\herdisc(\SS_1) \cdots \herdisc(\SS_k)}
      \lesssim
      \log(2m)^{k + \frac12},
    \end{multline*}
    where $m = |\SS_1|\cdots |\SS_k|$.
  \end{enumerate}
\end{theorem}
\begin{proof}
  Each of the bounds follows from Theorem~\ref{thm:ch7-fact-main}, and
  respectively, properties~\ref{pr:ch7-transpose}, \ref{pr:ch7-union}, and
  \ref{pr:ch7-tensor} of Lemma~\ref{lm:ch7-gamma2-props}. 
\end{proof}

\section{Applications}

In this section we outline several applications of
Theorem~\ref{thm:ch7-fact-main} to proving upper and lower bounds on
the discrepancy of set systems.

\subsection{The $\gamma_2$ Norm of Intervals}

We start with a set system that, at first sight appears to be quite
uninteresting: the set system $\II_n$ of prefix intervals  of $[n]$, i.e., sets
of the type $\{1, \ldots i\}$. The hereditary discrepancy of $\II_n$
is $1$: we just need to alternate the signs the elements we color to
achieve this bound. It turns out however, that the $\gamma_2$ norm of
the incidence matrix of $\II_n$ is on the order of $\log
n$. This fact, which we prove below, shows that the lower bound in 
Theorem~\ref{thm:ch7-fact-lb} is tight up to constants. It is also
useful in the applications we develop below.

We use $T_n$ for the incidence matrix of $\II_n$, which we
write as the lower triangular matrix
\[
T_n(i,j) =
\begin{cases}
  1 &i \ge j\\
  0 &i < j
\end{cases}.
\]

\begin{lemma}\label{lm:ch7-gamma2-Tn}
  For the matrix $T_n$ above, $\gamma_2(T_n) = \Theta(\log n)$. 
\end{lemma}
\begin{proof}
  The upper bound follows from Theorem~\ref{thm:ch7-fact-lb}, since
  \[
    \gamma_2(T_n) \lesssim (1 + \log n) \herdisc(T_n) = 1 + \log n.
  \]
  For the lower bound $\gamma_2(T_n) \gtrsim 1+\log n$, we use
  Lemma~\ref{lm:ch7-fact-dual}. In particular, setting $P = Q =
  \frac{1}{\sqrt{n}}$ in \eqref{eq:ch7-fact-dual}, we see that
  \(
    \gamma_2(T_n) \ge \frac1n \|T_n\|_{\tr}.
  \)
  To prove a lower bound on $\|T_n\|_{\tr}$, we need to estimate the
  singular values of $T_n$. This can be done directly by computing the
  eigenvectors of $T_n^T T_n$, which have a simple closed
  form. Instead, we will consider a matrix that is easier to
  diagonalize. Consider the matrix $S_n = 2T_n - J_n$, where $J_n$ is
  the $n\times n$ all-ones matrix. In other words, $S_n$ is the matrix
  we get from $T_n$ by replacing all entries equal to $0$ with
  $-1$. By the triangle inequality,
  \begin{equation}\label{eq:ch7-Tn-lb}
    \gamma_2(T_n) \ge \frac1n \|T_n\|_{\tr}
    \ge \frac{1}{2n} (\|S_n\|_{\tr} - \|J_n\|_{\tr})
    = \frac{1}{2n} \|S_n\|_{\tr} - \frac12,
  \end{equation}
  since $J_n$ has only one nonzero singular value equal to $n$. Notice
  that $S_n$ is a normal matrix, i.e., $S_n^T S_n = S_nS_n^T$, and,
  therefore, the singular values of $S_n$ are equal to the absolute
  values of the eigenvalues of $S_n$. To compute $\|S_n\|_{\tr}$, it
  then suffices to compute its eigenvalues. We can diagonalize $S_n$
  over the complex numbers, using a transformation similar to the
  discrete Fourier transform. Let $\omega \in \mathbb{C}$ be a primitive $2n$-th
  root of unity, and, for $k \in \{0, \ldots, n-1\}$,  define the
  vector $u_k$ by $u_k(i) = \omega^{(2k + 1)(i-1)}$. What motivates
  this choice is that $\omega^{(2k + 1)n} = -1$. It is easy to
  verify that $u_k$ and $u_\ell$ are orthogonal (with the standard
  inner product over $\mathbb{C}^n$) if $k \neq \ell$. Moreover,
  notice that
  \begin{align*}
    (S_nu_k)(i)
    &=
      (1 + \ldots + \omega^{(2k + 1)(i-1)})
      - (\omega^{(2k + 1)i} +\ldots +  \omega^{(2k + 1)(n-1)})\\
    &= \frac{\omega^{(2k + 1)i} - 1}{\omega^{2k + 1} - 1}
      - \frac{\omega^{(2k + 1)n} - \omega^{(2k + 1)i}}{\omega^{2k + 1} - 1}\\
    &= \frac{2\omega^{(2k + 1)i}}{\omega^{2k + 1}-1}
      = \frac{2\omega^{2k + 1}}{\omega^{2k + 1}-1}u_k(i).
  \end{align*}
  Therefore, $u_k$ is an eigenvector of $S_n$ with eigenvalue
  $\frac{2\omega^{2k + 1}}{\omega^{2k + 1}-1}$. Thus, the
  singular values of $S_n$ are $2/|\omega - 1|,
  \ldots, 2/|\omega^{2n-1} - 1|$. Using elementary
  trigonometry,
  \(
  |\omega^{2k + 1} - 1| = 2\sin\left(\frac{\pi(2k+1)}{2n}\right),
  \)
  and, therefore,
  \[
    \|S_n\|_{\tr} = \sum_{k = 0}^{n-1} \frac{1}{\sin\left(\frac{\pi(2k+1)}{2n}\right)}
    \ge
    \frac{2n}{\pi}\sum_{k = 0}^{n-1} \frac{1}{2k+1},
  \]
  as $\sin x \le x$ for $x \ge 0$.
  Plugging back into \eqref{eq:ch7-Tn-lb}, we get that
  \[
    \gamma_2(T_n) \ge \frac{1}{\pi}\sum_{k = 0}^{n-1} \frac{1}{2k+1}
    \gtrsim 1 + \log n. \qedhere
  \]
%  as we wanted to prove. 
\end{proof}

\begin{exercise}
  Using canonical intervals (see Section~\ref{sec:ch2-tusnady}), give a factorization
  showing that $\gamma_2(T_n) \lesssim 1 + \log n$. This gives a more
  direct proof of the upper bound in Lemma~\ref{lm:ch7-gamma2-Tn}.
\end{exercise}

\subsection{The Discrepancy of Boxes}

Let us now use what we have proved to show nearly tight bounds on the
discrepancy of boxes in $d$-dimensions, resolving
Problem~\ref{prob:tusnady} up to a fixed power of log.

Recall that
$\RR_d$ is the set of all axis aligned boxes in $\R^d$, i.e., sets of
the type $\rbox_{a,b} = \{u \in \R^d: a \le u \le b\}$ for $a,b \in
\R^d$. Recall also that $\disc(\RR_d, n)$ is the largest discrepancy
achieved by restricting $\RR_d$ to a set $P$ of $n$
points. Theorem~\ref{thm:ch7-fact-main}, Lemma~\ref{lm:ch7-gamma2-Tn},
and the properties of the $\gamma_2$ norm allow us to prove the
following bounds on $\disc(\RR_d,n)$.

\begin{theorem}\label{thm:ch7-tusnady}
  For any integer $d \ge 1$,
  \[
    \log(2n)^{d-1} \lesssim_d \disc(\RR_d,n) \lesssim_d \log(2n)^{d+\frac12},
  \]
  where the $\lesssim_d$ means that the constant
  depends on $d$.
\end{theorem}
\begin{proof}
  As we have done before with boxes, we will use the fact that we only
  need to consider anchored boxes, restricted to the grid
  $[n]^d$. Recall that $\AA_d$ is the set system of anchored boxes of the type $[0,b_1]\times\cdots \times [0,b_d]$ for $b \in \R^d$, and consider
  the set system $\SS \eqdef \AA_d|_{[n]^d}$ of subsets of the
  $n\times \cdots \times n$ grid induced by anchored boxes. As we observed in Section~\ref{sec:ch2-tusnady},
  \begin{equation}\label{eq:ch7-tusnady-grid}
   2^{-d}\disc(\RR_d,n) \le \herdisc(\SS) \le \disc(\RR_d,n^d).
 \end{equation}
  Thus, it suffices to bound $\herdisc(\SS)$ from
  both sides, and we do so by computing the $\gamma_2$ norm of the
  incidence matrix of $\SS$.

  Let then $\mat$ be the incidence matrix of $\SS$. The main
  observation is that $\SS = \II_n^{\times d}$, i.e., the product of
  $\II_n$ with itself $d$ times. Therefore,  $\mat =
  T_n^{\otimes d}$, i.e., the Kronecker
  product of $T_n$ with itself $d$
  times. Lemma~\ref{lm:ch7-gamma2-props} then shows that
  \(
  \gamma_2(\mat) = \gamma_2(T_n^{\otimes d}) = \gamma_2(T_n)^d,
  \)
  so, by Lemma~\ref{lm:ch7-gamma2-Tn}, 
  \begin{equation}\label{eq:ch7-gamma2-tusnady}
    \log(2n)^d \lesssim_d \gamma_2(\mat) \lesssim_d \log(2n)^d.
  \end{equation}
  Theorem~\ref{thm:ch7-fact-main} now implies
  \[
    \log(2n)^{d-1} \lesssim_d \herdisc(\SS) \lesssim_d \log(2n)^{d+\frac12}
  \]
  since $|\SS| \le n^d$. This implies the theorem via
  \eqref{eq:ch7-tusnady-grid}.
\end{proof}

Below we will see that the upper bound in
Theorem~\ref{thm:ch7-tusnady} can be improved further. The lower bound
is the best one known for Tusn\'ady's problem.

\subsection{Other Applications}

Let us note a couple of other discrepancy upper bounds that can be
proved using Lemma~\ref{lm:ch7-gamma2-props},
Theorem~\ref{thm:ch7-props-ss}, and
Theorem~\ref{thm:ch7-fact-main}. Neither of the two bounds below is
optimal, but they are both non-trivial and follow easily from the
results proved thus far.

\paragraph{Permutation families.} The set system induced by a single
permutation of $[n]$ clearly has hereditary discrepancy $1$, as we can list the
elements of $[n]$ (or a subset of it) in the order given by the permutation, and alternate
the signs of the coloring. Then, by property~\ref{pr:ch7-ss-union}.~of
Theorem~\ref{thm:ch7-props-ss}, we have that, for a set system $\SS$
induced by $k$ permutations of $[n]$, $\disc(\SS) \lesssim \sqrt{k} \log(kn)^{3/2}$.

\paragraph{Arithmetic Progressions.} Consider the system $\mathcal{AP}$
consisting of all arithmetic progressions restricted to $[n]$. A
classical result of Roth shows that $\disc(\mathcal{AP}) \gtrsim n^{1/4}$. We
can show an upper bound that matches Roth's lower bound up to
logarithmic factors. Let $\mathcal{AP}_d$ be the subset of $\mathcal{AP}$ consisting of
arithmetic progressions of difference $d$, i.e., sets of the form $a +
d, a+2d, \ldots, a + \ell d$ for some integers $a$ and $\ell$. Let
$\mathcal{AP}_{\ge s} \eqdef \bigcup_{d \ge s} \mathcal{AP}_{d}$ for a parameter $s$ we will
choose soon. Since
every arithmetic progression in $\mathcal{AP}_{\ge s}$ has difference at least
$s$, it can have size at most $\frac{n}{s}$, so, for the incidence
matrix $\mat_{\ge s}$ of $\mathcal{AP}_{\ge s}$, $\gamma_2(\mat_{\ge s}) \le
\sqrt{\frac{n}{s}}$. On the other hand, the hereditary discrepancy of
$\mathcal{AP}_d$ for any $d$ is at most $1$, since we can separately color each
congruence class $\mod d$, and alternate signs of the coloring within
a class. Therefore, by Theorem~\ref{thm:ch7-fact-main},
$\gamma_2(\mat_{d}) \lesssim \log(2n)$ for the incidence matrix
$\mat_d$ of $\mathcal{AP}_d$. Using property~\ref{pr:ch7-union} of
Lemma~\ref{lm:ch7-gamma2-props} we have, for the incidence matrix
$\mat$ of $\mathcal{AP}$,
\begin{align*}
  \gamma_2(\mat) &\le (\gamma_2(\mat_{\ge s})^2 +
  \gamma_2(\mat_{s-1})^2 + \cdots + \gamma_2(\mat_1)^2)^{1/2}\\
  &\lesssim \left(\frac{n}{s} + s\log(2n)^2\right)^{1/2}.
\end{align*}
Optimizing over $s$, we have $\gamma_2(\mat) \lesssim
n^{1/4}\sqrt{\log(2n)}$, and Theorem~\ref{thm:ch7-fact-main} gives us 
\(
\disc(\mathcal{AP}) \lesssim n^{1/4}\log(2n).
\)

\begin{exercise}
  Use canonical intervals to show that $\gamma_2(\mat)\lesssim
  n^{1/4}$ for the incidence matrix $\mat$ of the set system of
  arithmetic progressions restricted to $[n]$.
\end{exercise}
\section{Prefix Discrepancy}

Recall that in Lemma~\ref{lm:fact-ub-prefix} we showed that, for any
matrix $\mat$ with columns $u_1, \ldots, u_n \in \R^m$, the prefix discrepancy
\begin{equation}\label{eq:ch7-ub-prefix}
  \pdisc(\mat) \lesssim \gamma_2(\mat)\sqrt{\log(2mn)}.
\end{equation}
% where $\pdisc(\mat)$ is the minimum over $\col \in \{-1,+1\}^n$ of 
% \[
%   \pdisc(\mat,\col) = \max_{j = 1}^n \left\|\sum_{i = 1}^j \col(i)u_u\right\|_\infty.
% \]
Theorem~\ref{thm:ch7-fact-main} and this inequality show that prefix
discrepancy is never much bigger than the hereditary discrepancy,
since, for any $m\times n$ matrix $\mat$,
\begin{equation}
  \label{eq:ch7-pdisc-herdisc}
  \pdisc(\mat) \lesssim \log(mn)^{1.5} \herdisc(\mat).
\end{equation}

In addition to this connection between prefix and hereditary
discrepancy, the upper bound~\eqref{eq:ch7-ub-prefix} sometimes allows us to
prove tighter discrepancy upper bounds than
Theorem~\ref{thm:ch7-fact-main}. We give two such examples below, the
first giving the best known  upper bound for Tusn\'ady's problem.

\medskip
\noindent {\bf Axis-aligned boxes.} As we will see in the next theorem, the
bound \eqref{eq:ch7-pdisc-herdisc} allows us to improve the upper
bound in Theorem~\ref{thm:ch7-tusnady} by effectively getting one
dimension ``for free''.

\begin{theorem}\label{thm:ch7-tusnady-ub}
  For any integer $d \ge 1$,
  \(
    \disc(\RR_d,n) \lesssim_d \log(2n)^{d-\frac12}.
  \)
\end{theorem}
\begin{proof}
  As in Theorem~\ref{thm:ch7-tusnady}, we only consider anchored
  boxes, restricted to sets in $[n]^d$. Equation
  \eqref{eq:ch7-pdisc-herdisc} allows us to restrict the sets we
  consider even more. Let $P \subseteq [n]^d$ be some $n$-point set,
  and let $\SS$ be the set system of all subsets of
  $P$ of the form $R \cap P$ where
  \(
  R = [0,b(1)] \times\cdots \times [0,b({d-1})]\times \R,
  \)
  for $b\in \R^{d-1}$. That is, $\SS$ consists of sets
  induced by anchored boxes that extend to infinity in the last
  coordinate. Equivalently, let $\pi_{d-1}$ be the orthogonal
  projection onto the first $d-1$ coordinates of $\R^d$; then $\SS$ is
  the set system consisting of all sets of the type $\{p: \pi_{d-1}(p)
  \in R\}$ where $R \in \AA_{d-1}$ is a $(d-1)$-dimensional anchored
  box. Since $\SS$ is isomorphic to $\AA_{d-1}|_{\pi_{d-1}(P)}$, equation \eqref{eq:ch7-gamma2-tusnady} in the proof
  of Theorem~\ref{thm:ch7-tusnady} gives us that the incidence matrix
  $\mat$ of $\SS$ satisfies $\gamma_2(\mat)\lesssim_d
  \log(2n)^{d-1}$.


  Let us order $P$ as $p_1, \ldots, p_n$ so that $p_1(d) \le \ldots
  \le p_n(d)$, and also order the columns of the incidence matrix $\mat$ of $\SS$
  so that the $i$-th column is the indicator vector of sets containing $p_i$.
  The key observation we make is that, for any anchored box $R \in
  \AA_d$, $R \cap P = S \cap \{p_1, \ldots, p_i\}$ for some $S \in
  \SS$ and $i \in [n]$. In particular, we can take $p_i$ to be the
  point with the largest $p_i(d)$ coordinate contained in $R$, and then
  take $S = \{p \in P: \pi_{d-1}(p) \in \pi_{d-1}(R)\}$. Therefore, for any
  coloring $\col:P \to \{-1,+1\}$, $|\col(R \cap P)|$ is bounded by
  $\pdisc(\mat, \col)$. The theorem now follows from
  \eqref{eq:ch7-ub-prefix} since
  \[
    \disc(\AA_d|_P) \le \pdisc(\mat)
    \lesssim_d \sqrt{\log(2n)} \gamma_2(\mat)
    \lesssim_d \log(2n)^{d-\frac12},
  \]
  where we used the fact that $\AA_d|_P$ consists of at most $n^d$
  distinct sets, as well as the estimate
  \(\gamma_2(\mat)\lesssim_d  \log(2n)^{d-1}\) above.
\end{proof}

\paragraph{Arithmetic Progressions.} Recall that $\mathcal{AP}$ is the
set system of arithmetic progressions restricted to $[n]$. We can use
equation~\eqref{eq:ch7-pdisc-herdisc} to improve our upper bound on
its discrepancy to $\disc(\mathcal{AP}) \lesssim n^{1/4}\sqrt{\log
  n}$. To see this, let us define $\SS$ to be the set of
\emph{maximal} arithmetic progressions, \ie, progressions inside $[n]$
that cannot be extended further to the left or to the right. We claim
that the incidence matrix $\mat$ of $\SS$ satisfies $\gamma_2(\mat)
\lesssim n^{1/4}$. To see this, notice that
\(
\mat =
\begin{pmatrix}
  \mat_{\ge \sqrt{n}}\\
  \mat_{< \sqrt{n}}
\end{pmatrix}
\), where $\mat_{\ge \sqrt{n}}$ and $\mat_{< \sqrt{n}}$ are the incidence matrix of maximal
arithmetic progressions of difference at least $\sqrt{n}$, and at most $\sqrt{n}$  respectively. 

Recall that each arithmetic progression in $[n]$ of difference $d$ has
cardinality at most $n/d$, and all maximal arithmetic
progressions of the same difference are disjoint. Thus, it follows that $\gamma_2(\mat_{\ge \sqrt{n}}) \le n^{1/4}$ as each of its rows has at
most $\sqrt{n}$ ones, and $\gamma_2(\mat_{<\sqrt{n}}) \le n^{1/4}$
as each of its columns has at most $\sqrt{n}$ ones, one for each
possible difference. So  property~\ref{pr:ch7-union} of
Lemma~\ref{lm:ch7-gamma2-props} gives us that $\gamma_2(\mat)\le
\sqrt{2} n^{1/4}$.  Now equation
\eqref{eq:ch7-ub-prefix} gives us that
\(
\pdisc(\mat) \lesssim n^{1/4} \sqrt{\log n},
\)
where the colums of $\mat$ are ordered so that the $i$-th column is
the indicator of the sets in $\SS$ containing $i$. This means that
there exists a coloring $\col:[n]\to \{-1,+1\}$ such that $|\col(S
\cap [j])| \lesssim n^{1/4} \sqrt{\log n}$ for each $j \in [n]$ and $S
\in \SS$. Since any arithmetic progression inside $[n]$ can be written
as $(S \cap [k]) \setminus (S \cap [j])$, we also have
$\disc(\mathcal{AP},\col) \lesssim n^{1/4}\sqrt{\log n}$.

\section{Bibliographic Notes}

The computational complexity of discrepancy was first investigated by
Charikar, Newman and Nikolov, who proved
Theorem~\ref{thm:ch7-disc-hard}, and analogous results for set systems
of bounded degree, and for set systems of bounded VC
dimension~\cite{dischard}. It was shown by Austrin, Guruswami, and
H{\aa}stad that hereditary discrepancy is $\mathsf{NP}$-hard to
approximate within factor $2-\varepsilon$ for any
$\varepsilon > 0$~\cite{AustrinGH13}. 

The $\gamma_2$ norm was first introduced in Banach space theory as a
way to measure how well a bounded linear operator $u:X \to Y$ between
Banach spaces $X$ and $Y$ factors through a Hilbert space. In
particular, $\gamma_2(u)$ is the infimum of the product of operator
norms $\|v\|\|w\|$ over all operators $v:X \to H$ and $w:H \to Y$ such
that $wv = u$ and $H$ is some Hilbert space. Analogously, if we
replace the Hilbert space with an $L_p$ space over some measure space
$(\Omega,\mu)$, we get the $\gamma_p$ norm. The $\gamma_2$ norm we
define here for matrices corresponds to the special case when $X =
\ell_1^n$ and $Y = \ell_\infty^m$.
Factorization of linear operators through
$L_p$ spaces was studied systematically by Kwapien~\cite{Kwapien72},
but has its origins already in the work of
Grothendieck~\cite{Grothendieck52}. In particular, the dual formulation
\eqref{eq:ch7-fact-dual} can be derived from the results of
Grothendieck (up to a factor of 2) and Kwapien; see also the appendix
of Pisier's survey on Grothendieck's inequality~\cite{Pisier-grothendieck}. The book by
Tomczak-Jaegermann surveys the basic facts and many
applications of the $\gamma_p$ norms in Banach space theory~\cite{Tomczak-Jaegermann}.

The connection between hereditary discrepancy and the $\gamma_2$ norm
is due to Matou\v{s}ek, Nikolov, and Talwar, who also proved
Theorems~\ref{thm:ch7-fact-main}~and~\ref{thm:ch7-tusnady}~\cite{MNT}. The
semidefinite formulation in Lemma~\ref{lm:ch7-fact-sdp}, and the
computability and duality results in Lemma~\ref{lm:ch7-fact-dual} are
due to Lee, Shraibman, and \v{S}palek~\cite{LeeSS08}. 

The $\gamma_2$ norm has also found other applications, \eg, to
communication complexity and learning
theory~\cite{LinialMSS07-signmatrices} (see also the
book~\cite{LeeS09}), and differential privacy~\cite{NTZ}. It is also
implicit in the work of Larsen on lower bounds for oblivious range
searching data structure~\cite{disc-larsen}, who also discovered the discrepancy upper
bounds stated in Lemma~\ref{lm:fact-ub} and
Corollary~\ref{cor:ch7-fact-ub}.

The determinant lower bound in Lemma~\ref{lm:ch7-detlb} was first
shown by Lov\'asz, Spencer, and Vesztergombi~\cite{LSV}. They also asked
if hereditary discrepancy can be bounded from above by a function of
the determinant lower bound. This was disproved by Hoffman, and his
counterexample can be found in Chapter~4 of Matou\v{s}ek's
book~\cite{matousek2010geometric}. Nevertheless, it is possible to
show that $\herdisc(\mat) \lesssim
\detlb(\mat)\sqrt{\log(2m)\log(2n)}$. The first result of this type
was proved by Matou\v{s}ek~\cite{Matousek-detlb}, and his (loose)
bound was strengthened by Jiang and
Reis~\cite{JiangReis22}.  Li and Nikolov showed this bound is
tight~\cite{LiNikolov24}. 
The volume
lower bound in Remark~\ref{rem:vollb} was observed by
Banaszczyk~\cite{Bana93}. It was further investigated by Dadush, Nikolov,
Talwar, and Tomczak-Jaegermann, who showed it matches the hereditary
discrepancy up to logarithmic factors for any finite dimensional norm,
and matches a hereditary notion of partial coloring discrepancy up to
constant factors~\cite{anynorm}.

The lower bound on $\gamma_2(T_n)$ in Lemma~\ref{lm:ch7-gamma2-Tn} is due to Mathias~\cite{Mathias}. He also showed that $\gamma_2(S_n) = \frac{1}{n} \|S_n\|_{\tr}$, and, therefore, $\gamma_2(T_n) \le \frac{1}{2n}\|S_n\| + \frac12$. Together with \eqref{eq:ch7-Tn-lb}, this gives us the value of $\gamma_2(T_n)$ up an additive approximation of $1$. Determining the exact value of $\gamma_2(T_n)$ as a function of $n$, and giving an optimal factorization remain intriguing open problems, with applications in differential privacy~\cite{DP-FTRL}.

The improved upper bound on the discrepancy of boxes in
Theorem~\ref{thm:ch7-tusnady-ub} is due to
Nikolov~\cite{tusnady-ub}. Essentially the same proof, together with a
decomposition theorem of Matou\v{s}ek~\cite{Matousek99-boxes}, gives similar upper bounds for
more general set systems induced by polytopes with prescribed facets. 
Moreover, the lower bound in Theorem~\ref{thm:ch7-tusnady}  also
extends to these more general set systems, under some technical conditions. These
results are also due to Nikolov~\cite{tusnady-ub}. 

The discrepancy of arithmetic progressions is bounded by $O(n^{1/4})$,
and this was shown by Matou\v{s}ek and Spencer using the partial
coloring method~\cite{MatousekSpencer-ap}: see the notes to
Chapter~\ref{ch:partial} for more on the discrepancy of arithmetic
progressions. 




%%% Local Variables:
%%% mode: latex
%%% TeX-master: "book"
%%% End:
